\chapter{CONCEITOS GERAIS E REVISÃO DA LITERATURA}

%Lembre-se que os capítulo, as sessões e sub-sessões são determinadas por si para adequar-se ao seu trabalho.

Neste capítulo deve ser proporcionado o estado da arte / referencial teórico sobre o tema a que se refere o estudo. Um bom pesquisador não deve repetir trabalhos já concluídos ou que já estão em andamento. Por isso esta sessão é onde o autor demonstra até onde vai a pesquisa atual no campo de estudos em questão e estabelece as bases sobre as quais desenvolverá o estudo proposto. 

\section{Conceitos Gerais}

\subsection{Exemplos de Objetos do Texto}

\subsubsection{Citações}
As citações em textos tal como \citeonline[p. 1]{examplephdthesis}, utilize o comando \textit{"citeonline"}, entre colchetes (opcional) está a referencia da página do trabalho, exemplo sem colchetes \citeonline{examplephdthesis}.

Citação em final de parágrafo utilize o comando cite. \cite{examplephdthesis}

O argumento (entre colchetes) também é opcional. \cite[p. 1]{examplephdthesis}

\subsubsection{Estruturas Matemáticas}
A seguir são mostrados alguns exemplos de como deve-se utilizar ambiente teorema, axioma, corolário, lema e demonstrações assim como inserir as figuras, citações diretas e tabelas. A Figura 01 mostra um exemplo de como inserir uma figura no texto. A Tabela 01 mostra o exemplo de como uma tabela deve ser inserida.

Para axiomas:
\begin{axm}[**Opcional o nome do axioma]
    Coisas que são iguais a uma mesma coisa, são iguais entre si.
\end{axm}

\begin{defi}[**Opicional o nome da definição]
A definição define o que deve ser definido.
\end{defi}

Para teoremas:
\begin{thm}[**Opcional o nome do teorema]
    \label{thm}
    Sejam $a, b$ os catetos e $c$ a hipotenusa de um triangulo retângulo, então $$ a^2+b^2=c^2$$
\end{thm}

Com sua demonstração
\begin{proof}
    Demonstração do teorema \ref{thm}, é trivial.
    
    \cqd % comando para inserir quadrado preto no fim da demonstração, se retirá-lo não aparecerá o quadrado preto.
\end{proof}

Para proposição:
\begin{prop}
Proposição é um termo usado em lógica para descrever o conteúdo de asserções. Uma asserção é um conteúdo que pode ser tomado como verdadeiro ou falso. Asserções são abstrações de sentenças não linguísticas que a constituem. A natureza das proposições é altamente controversa entre filósofos, muitos dos quais são céticos sobre a existência de proposições. Muitos lógicos preferem evitar o uso do termo proposição em favor de usar sentença
\end{prop}

Para corolários:
\begin{cor}[**Opcional o nome do corolário]
    \label{cor}
    Um corolário (do latim tardio corollarĭum) é uma afirmação deduzida de uma verdade já demonstrada. Assim como proposição resultante de uma verdade.
\end{cor}

Para lemas:
\begin{lem}[**Opcional o nome do lema]
    \label{lem}
    A palavra "lema" vem do grego $\lambda\acute{\eta}\mu\mu\alpha$ (lémma), no sentido de 'proposição'.
\end{lem}

\subsubsection{Equações}

Exemplo de equação com simbolo de integral
\begin{equation}
 \dint_{x=3}^{x=4} \dfrac{x^2+3}{x-a} \mbox{d}x
\end{equation}


Exemplo de como inserir equação por partes.
\begin{equation}
 f(x) = \left\{ 
 \begin{array}{ll}
  x^2+3x+1 & , x \geq 0 \\
  \ln |x^3|, x < 0
 \end{array}\right.
\end{equation}

Exemplo de contrução de matriz $3\times3$
\begin{equation}
  \left[
 \begin{array}{ccc}
  2 & 3 & 4 \\
  3 & 4 & 5 \\
  6 & 7 & 8 
 \end{array}
  \right]
\end{equation}

\subsection{Estruturas Textuais}
Para citação direta, teremos:
\begin{citacao}
	A Matemática é a chave do portão e as ciências. A falta de atenção às obras matemáticas prejudica todos os conhecimentos, uma vez que ele é ignorante de não poder conhecer as outras ciências ou as coisas deste mundo. \cite{exampleRogerCitacao}
\end{citacao}

\subsubsection{Figuras}
Para figuras:

\begin{figure}[h]
\centering
\caption{Exemplo de como inserir Figura}
\includegraphics[scale=0.75]{imagens/logo_latex.eps}
\label{fig:exemplo}
\fonte{acervo do texto-exemplo} 
\end{figure}

Referenciando a figura \ref{fig:exemplo} como se segue. Para inserir duas figuras será da seguinte forma:

\begin{figure}[h]
    \centering
    \caption{Imagens utilizadas}
    \begin{subfigure}[b]{0.3\textwidth}
        \centering
        \caption{Passarinho}
        \includegraphics[scale=0.5]{imagens/logo_latex.eps}
        \label{fig:pas}
    \end{subfigure}
    \qquad
    \begin{subfigure}[b]{0.3\textwidth}
        \centering
        \caption{Brasão Universidade}
        \includegraphics[scale=1]{imagens/logoifma.eps}
        \label{fig:uni}
    \end{subfigure}
    \fonte{acervo do texto-exemplo} 
    \label{fig:my_label}
\end{figure}

Podendo referenciar \ref{fig:pas} e \ref{fig:uni} como agora. Todas as imagens devem ser referenciadas, mesmo que sejam de autoria propria. Elas devem estar em extensão .EPS, para uma melhor visualização.

\subsubsection{Tabelas}

Para tabelas:
\begin{table}[h]
  \centering
  \caption{ Modelo de como as tabelas devem ser inseridas no texto }
  \begin{tabular}{|c|c|c|c|}
    \hline 
    \textbf {Índice} & \textbf{Coluna 01} &\textbf{ Coluna 02} & \textbf{Coluna 03} \\ \hline \hline
    Linha 01 & & & \\ \hline
    Linha 02 & & & \\ \hline                         
  \end{tabular}
  \fonte{Criação para o texto-exemplo} 
  \label{tab:modelo}
\end{table}


